\section{Chapitre~\ref{sec:livingmachine}. Une machine faite de neurones vivants}

Le comportement des cultures sur matrices d'électrodes est mesuré dans des expériences directes d'enregistrement et de stimulation; le mécanisme des salves s'appuie sur le modèle d'épuisement synaptique et sur des expériences en microcultures, et les affirmations sur la dopamine en boîte de culture sur des expériences isolées, dont certaines ne sont, de l'aveu même des auteurs, qu'une démonstration.

{\footnotesize
\setlength{\tabcolsep}{3pt}\renewcommand{\arraystretch}{1.15}
\begin{longtable}{>{\raggedright}p{0.5cm}>{\raggedright}p{4.0cm}>{\raggedright}p{5.6cm}>{\raggedright}p{2.3cm}>{\raggedright\arraybackslash}p{1.7cm}}
\toprule
N\textsuperscript{o} & Affirmation & Expérience et ce qui est mesuré & Source & Statut \\
\midrule\endfirsthead
\toprule
N\textsuperscript{o} & Affirmation & Expérience et ce qui est mesuré & Source & Statut \\
\midrule\endhead
1 & Culture sur puce: surtout des \nt{GLUT}, une part plus faible de \nt{GABA} qui dépend de la préparation; environ $6\%$ dans les cultures d'hippocampe & Un réseau de milliers de neurones sur matrice d'électrodes est une expérience standard (ligne~3). La part des neurones \nt{GABA} dans les cultures d'hippocampe a été comptée par marquage du GABA: environ $6\%$ des neurones. Pour les cultures de cortex, nous ne donnons pas de chiffre vérifié & \cite{Benson1994} & mesuré \\
2 & Organoïdes: tissu nerveux tridimensionnel d'environ un millimètre, issu de cellules souches humaines & À partir de cellules souches, on a cultivé des organoïdes tridimensionnels comportant plusieurs régions cérébrales, dont un cortex avec des couches de progéniteurs et des neurones matures & \cite{Lancaster2013} & mesuré \\
3 & Sans entrée, la culture s'embrase en salves espacées d'une seconde à plusieurs minutes, selon la culture & $58$ cultures de cortex, d'une densité de $3\,000$ à $50\,000$ neurones, ont été enregistrées pendant cinq semaines ($963$ enregistrements d'une demi-heure): après deux semaines, les salves de tout le réseau dominent dans presque toutes les cultures; les intervalles entre salves vont de $1$ à $300$~s et dépendent fortement de la préparation & \cite{Wagenaar2006} & mesuré \\
4 & Une salve s'achève par épuisement du neurotransmetteur, avec un intervalle $T_{\text{salve}}\approx\tau_{\text{rec}}\ln\frac{1}{1-x^*}$~\eqref{eq:burstinterval}; $2$--$4$~s est l'estimation du modèle avec $\tau_{\text{rec}}=0.8$~s, la récupération mesurée est plus lente & La formule est établie dans le chapitre. Modèle: un réseau à synapses dépressives produit spontanément des salves de tout le réseau. Expérience sur des microcultures de $4$--$30$ neurones: la durée d'une salve dépend du temps écoulé depuis la précédente, et l'épuisement des vésicules de neurotransmetteur supprime les salves; conclusion des auteurs: la salve est limitée par la réserve de neurotransmetteur. La réserve s'y reconstitue en quelques dizaines de secondes, ce qui, avec la même formule, donne des intervalles de dizaines de secondes à plusieurs minutes & démonstration dans le chapitre; \cite{Tsodyks2000,CohenSegal2011,Tsodyks1997} & théorie \\
5 & \og Le professeur qui se tait\fg{}: la stimulation est coupée quand la réponse voulue apparaît, et la réponse se fixe & La culture était stimulée à $0.3$--$1$~Hz jusqu'à ce que la réponse choisie apparaisse $50\pm10$~ms après le stimulus, puis la stimulation était coupée pendant $5$~min; après quelques cycles, la réponse voulue était évoquée d'emblée; des courbes d'apprentissage ont été tracées & \cite{Shahaf2001} & mesuré \\
6 & La culture joue au tennis; hausse significative des séries en une session seulement avec rétroaction complète & Détails au chapitre~\ref{sec:dishbrain} et dans son annexe: avec du silence à la place du stimulus, la culture joue aussi mieux en fin de session que sans rétroaction, mais avec un effet moindre; l'apprentissage d'une session à l'autre est instable & \cite{Kagan2022} & mesuré \\
7 & Que la culture apprenne à prédire son entrée reste une hypothèse; les grands mots sur sa \og sensibilité\fg{} sont contestés & Le principe de l'énergie libre est une proposition théorique; aucune expérience directe ne le distingue d'explications plus simples. Dans un article de réponse, trente chercheurs ont souligné qu'un changement de comportement dans la boucle ne démontre ni intelligence ni sensations & \cite{Friston2010,Balci2023} & hypothèse \\
8 & La culture dirige un corps virtuel vers une cible avec un motif de stimuli bien choisi & Un programme choisissait lui-même les motifs de stimulation d'entraînement selon le résultat courant; en quelques dizaines de minutes, le réseau atteignait les états visés, et la plasticité durait $80$~min après $2$~h d'entraînement. Les mêmes stimuli, appliqués sans lien avec le comportement, restaient sans effet & \cite{Bakkum2008} & mesuré \\
9 & Réservoir: seule la lecture linéaire $y=W_{\text{sort}}\mathbf r$~\eqref{eq:reservoir} est entraînée & Théorie du calcul par réservoir: tout système non linéaire à mémoire évanescente, plus une sortie linéaire entraînée & \cite{Maass2002,JaegerHaas2004} & théorie \\
10 & Un organoïde utilisé comme réservoir distingue les locuteurs avec une précision d'environ $78\%$ (selon les auteurs); la lecture apprend par l'erreur, et les connexions de l'organoïde changent sans professeur & Des sons de parole de plusieurs locuteurs étaient injectés à l'organoïde sous forme de motifs de courant sur une matrice d'électrodes; la lecture entraînée distinguait le locuteur avec une précision d'environ $78\%$, selon les auteurs. Ceux-ci rapportent aussi que la connectivité fonctionnelle de l'organoïde changeait pendant l'entraînement, et parlent d'apprentissage non supervisé dans l'organoïde lui-même & \cite{Cai2023} & mesuré \\
11 & Un million de neurones dépense environ $10^{-4}$~W en impulsions~\eqref{eq:dishpower} & Estimation: le coût d'une impulsion, $10^{-10}$~J, tiré du budget énergétique du cortex~\eqref{eq:energy}, multiplié par le nombre d'impulsions; il n'existe pas de mesure directe de la puissance d'une culture & démonstration dans le chapitre; \cite{Attwell2001} & théorie \\
12 & Dopamine par éclair de lumière: $1$~mM de dopamine en cage, $240$ répétitions, le nombre d'impulsions évoquées augmentait & Sur une plateforme d'organoïdes à distance, la dopamine enfermée dans une cage photosensible ($1$~mM) était libérée par ultraviolet ($365$~nm, $0.8$~s) quand le stimulus provoquait plus d'impulsions; en $240$ étapes (environ $5$~h), le nombre d'impulsions a globalement augmenté. Les auteurs écrivent qu'une seule expérience ne permet pas d'établir le lien avec la dopamine; il n'y a pas de contrôles & \cite{Jordan2024} & hypothèse \\
13 & La dopamine de cellules vivantes modifie le poids d'un transistor organique de façon durable et réversible & Des cellules sécrétant de la dopamine ont été cultivées au-dessus d'un composant neuromorphique organique; l'oxydation de la dopamine sur l'électrode modifie la conductance du canal; un changement durable du poids et son retour ont été démontrés & \cite{Keene2020} & mesuré \\
14 & Des organoïdes contenant des neurones \nt{DOPA} fonctionnels existent; leur dopamine n'a pas été utilisée comme professeur & Dans des organoïdes issus de cellules souches humaines, on a trouvé des neurones \nt{DOPA} matures électriquement actifs et une production de dopamine. Nous n'avons trouvé dans la littérature aucune expérience où cette dopamine enseigne à un autre réseau & \cite{Jo2016} & mesuré \\
15 & Un organoïde tient un pendule en équilibre: les stimuli d'apprentissage choisis par le programme font mieux que le hasard, sans consolidation, via les synapses \nt{GLUT} & Organoïdes de cortex de souris en boucle fermée sur la tâche du \og pendule inversé sur chariot\fg{}: pour la plupart des organoïdes, les signaux choisis par apprentissage par renforcement ont donné un meilleur résultat que des signaux aléatoires ou l'absence de signal; après $45$~min de repos, l'amélioration ne subsistait pas; le blocage des récepteurs AMPA et NMDA la supprimait & \cite{Robbins2026} & mesuré \\
16 & Sans registres de contrôle, le réseau du banc ne peut pas décider lui-même de ce qui compte comme un succès (proposition~\ref{prop:livingmachine}) & Dans toutes les expériences des lignes 5--15, le signal d'apprentissage est fixé par l'expérimentateur ou par un programme; aucune expérience ne montre une culture qui élaborerait elle-même un critère de succès. C'est une conclusion tirée d'une absence, et non une expérience & --- & hypothèse \\
\bottomrule
\end{longtable}}
